\documentclass[10pt,a4paper]{article}
\usepackage[utf8]{inputenc}
\usepackage[spanish]{babel}
\usepackage{geometry}
\geometry{margin=2.5cm}
\usepackage{booktabs}
\usepackage{amsmath}
\usepackage{listings}
\usepackage{xcolor}
\usepackage{hyperref}
\usepackage{float}

\definecolor{codebg}{rgb}{0.95,0.95,0.95}
\lstdefinestyle{cpp}{
    language=C++,
    basicstyle=\ttfamily\small,
    backgroundcolor=\color{codebg},
    frame=single,
    breaklines=true,
    numbers=left,
    numberstyle=\tiny,
    keywordstyle=\color{blue},
    commentstyle=\color{green!50!black},
    stringstyle=\color{red},
    tabsize=4
}
\lstdefinestyle{python}{
    language=Python,
    basicstyle=\ttfamily\small,
    backgroundcolor=\color{codebg},
    frame=single,
    breaklines=true,
    numbers=left,
    numberstyle=\tiny,
    keywordstyle=\color{blue},
    commentstyle=\color{green!50!black},
    stringstyle=\color{red},
    tabsize=4
}

\lstset{style=python}

\title{\textbf{Informe Grupal: Prueba(Parte II) - Desarrollo grupal de estrategias}}
\author{Grupo -- INF8090 Computaci\'on Paralela y Distribuida}
\date{Mayo 2026}


\begin{document}
\maketitle

\section*{Datos del Grupo}
\begin{table}[H]
\centering
\begin{tabular}{p{6cm}p{7cm}}
\toprule
\textbf{Integrante} & \textbf{Correo / Identificaci\'on} \\
\midrule
Elliot Acevedo Zuñiga & eacevedoz@utem.cl\\
Benjamin Aballay Palma 2 & baballay@utem.cl\\ 
Felipe Martinez Gonzalez 3 & fmartinezgo@utem.cl\\
\bottomrule
\end{tabular}
\end{table}
\vspace{0.5cm}
\noindent\textbf{Curso:} INF8090 -- Computaci\'on Paralela y Distribuida\\
\textbf{Evaluaci\'on:} Prueba 1 -- Parte II Grupal\\
\textbf{Fecha:} Mayo 2026
\newpage

% ============================================================
\section{Ejercicio 1: Normalizaci\'on Z-Score con OpenMP en C++}
% ============================================================

\subsection{Diagn\'ostico y Selecci\'on de Estrategia}

Durante la evaluaci\'on individual (Parte I), no se dej\'o registro de una
estrategia preliminar escrita para este ejercicio. Por lo tanto, al iniciar la
Parte II, el grupo parti\'o desde cero para evaluar c\'omo abordar el c\'alculo
de la media, varianza y normalizaci\'on Z-Score sobre 6.4 GB de datos.

Al discutir las alternativas como grupo, nos planteamos inicialmente un enfoque
ingenuo: paralelizar el c\'alculo directamente usando \texttt{\#pragma omp critical}
o \texttt{atomic} para ir sumando la media y varianza global a medida que se
le\'ia cada fila.

Sin embargo, descartamos esa opci\'on y llegamos a la conclusi\'on de que la
mejor forma de hacerlo era mediante reducciones y variables locales (Estrategia
de Memoria Compartida Local) por las siguientes razones t\'ecnicas:

\begin{itemize}
    \item \textbf{Control de Dependencias y Sincronizaci\'on:} Hacer un bloqueo
        constante (\texttt{critical}) por cada celda v\'alida obligar\'ia a los
        hilos a esperar turno millones de veces, serializando el c\'odigo y
        destruyendo el rendimiento. En su lugar, decidimos usar la directiva
        \texttt{reduction(+:...)} de OpenMP para las sumas globales (media y
        varianza). OpenMP maneja los acumuladores privados por hilo en bajo
        nivel y los suma al final de forma ultra eficiente.
    \item \textbf{Manejo de Localidad y Memoria:} Como este es un problema
        fuertemente limitado por el ancho de banda de la memoria
        (\emph{memory-bound}), estructuramos el c\'odigo repartiendo las filas
        con un \emph{scheduling} est\'atico (\texttt{\#pragma omp parallel for}).
        Esto permite que los n\'ucleos de la CPU trabajen sin el problema de
        \emph{false sharing}, aprovechando la cach\'e local antes de sincronizar
        resultados.
\end{itemize}

\subsection{Diagn\'ostico del Problema}

Se requiere normalizar una matriz de punto flotante de dimensi\'on
$50\,000\,000 \times 16$ ($\sim$6.4 GB en precisi\'on doble) que contiene
valores faltantes (NaN) en aproximadamente el 5\% de sus celdas. La
normalizaci\'on Z-Score se define como:

\[
z_{ij} = \frac{x_{ij} - \mu_j}{\sigma_j}
\qquad
\mu_j = \frac{1}{N_j}\sum_{i \notin \text{NaN}} x_{ij}
\qquad
\sigma_j = \sqrt{\frac{1}{N_j}\sum_{i \notin \text{NaN}} (x_{ij} - \mu_j)^2}
\]

donde $\mu_j$ y $\sigma_j$ se calculan ignorando los NaN, y los NaN se
preservan en la salida. Adem\'as, se cuentan los valores at\'ipicos
(\emph{outliers}) definidos como $|z_{ij}| > 3$.

El problema presenta tres desaf\'ios fundamentales:
\begin{enumerate}
    \item \textbf{Volumen de datos}: 6.4 GB excede la cach\'e de cualquier CPU
        moderna, por lo que el cuello de botella es el ancho de banda de la
        memoria principal (DRAM), no la capacidad de c\'omputo.
    \item \textbf{Manejo de NaN}: Cada celda debe verificarse con
        \texttt{isnan()} antes de contribuir a los estad\'isticos o a la
        normalizaci\'on.
    \item \textbf{Tres fases de recorrido}: A diferencia de la normalizaci\'on
        min-max (que requiere una pasada de min/max y otra de normalizaci\'on),
        el Z-Score exige tres pasadas completas: media, varianza y
        normalizaci\'on, incrementando el tr\'afico de memoria en $\sim$33\%.
\end{enumerate}

\subsection{Datos}

Se gener\'o un dataset sint\'etico de $50\,000\,000 \times 16$ a partir de
coordenadas reales de viajes de Uber (archivos CSV p\'ublicos con $>2$ millones
de registros). Cada fila representa un viaje con 16 atributos num\'ericos:
coordenadas de recogida y destino, distancia, tarifa, propina, peajes,
velocidad, tiempo en tr\'afico, duraci\'on, multiplicador \emph{surge},
cantidad de pasajeros y calificaci\'on del conductor. Se inyect\'o 5\% de NaN
aleatorio. El archivo se guard\'o en formato binario plano (\texttt{double[]})
para lectura directa desde C++.

\subsection{Estrategia de Soluci\'on}

Se dise\~n\'o un pipeline de tres fases, cada una implementada con private
arrays por hilo para evitar \emph{false sharing} y una secci\'on
\texttt{critical} para combinar al final:

\begin{enumerate}
    \item \textbf{Fase 1 (Media y conteo v\'alido)}: Recorrido de lectura con
        arreglos privados \texttt{local\_sums[NCOLS]} y
        \texttt{local\_vals[NCOLS]}. Cada hilo acumula sumas y conteos en sus
        arreglos locales. Al finalizar, se combinan en una secci\'on
        \texttt{critical}. Luego $\mu_j = \text{sum}_j / \text{count}_j$.
    \item \textbf{Fase 2 (Varianza)}: Segundo recorrido de lectura similar,
        acumulando $\sum (x_{ij} - \mu_j)^2$ en arreglos privados
        \texttt{local\_sum\_sq[NCOLS]} y combinando con \texttt{critical}.
        Luego $\sigma_j = \sqrt{\text{sum\_sq}_j / \text{count}_j}$.
    \item \textbf{Fase 3 (Normalizaci\'on Z-Score y detecci\'on de at\'ipicos)}:
        Tercer recorrido con \texttt{\#pragma omp parallel for
        reduction(+:outliers)}. Cada celda v\'alida se transforma a
        $z = (x - \mu_j) / \sigma_j$; se cuenta como at\'ipico si
        $|z| > 3$. Los NaN se preservan.
\end{enumerate}

\subsubsection{Justificaci\'on T\'ecnica de OpenMP}

Se seleccion\'o OpenMP sobre MPI o hilos POSIX por las siguientes razones:

\begin{itemize}
    \item \textbf{Memoria compartida}: Toda la matriz cabe en la RAM del sistema
        (15 GB disponibles), por lo que no se requiere distribuci\'on entre nodos.
    \item \textbf{Directivas de alto nivel}: \texttt{parallel for} y
        \texttt{critical} permiten paralelizar con m\'inima intervenci\'on
        manual sobre sincronizaci\'on.
    \item \textbf{Granularidad gruesa}: El trabajo por fila es homog\'eneo
        (cada fila procesa 16 doubles), por lo que \texttt{schedule(static)}
        por defecto balancea bien la carga.
    \item \textbf{Escalabilidad esperada}: Al ser un problema
        \emph{memory-bound}, se anticipa que el speedup est\'e limitado por el
        ancho de banda de memoria, no por el n\'umero de n\'ucleos.
\end{itemize}

\subsubsection{Uso de reduction y arreglos privados}

Para evitar \emph{false sharing} y condiciones de carrera, las Fases 1 y 2
utilizan arreglos privados por hilo (\texttt{local\_sums}, \texttt{local\_vals},
\texttt{local\_sum\_sq}) en una regi\'on \texttt{parallel}, combinando al final
con una secci\'on \texttt{critical} (solo 16 sumas por hilo, overhead
despreciable). La Fase 3 usa \texttt{reduction(+:outliers)} que crea copias
privadas del contador y las combina al finalizar el bucle.

\subsection{Benchmark y Resultados}

\subsubsection{Registro del Entorno de Ejecuci\'on}

\begin{itemize}
    \item \textbf{Sistema operativo}: Linux 7.0.10-1-cachyos (Arch Linux)
    \item \textbf{CPU}: AMD Ryzen 7 5700G, 8 n\'ucleos f\'isicos, 16 hilos l\'ogicos
    \item \textbf{RAM}: 15 GB total, $\sim$11 GB disponibles
    \item \textbf{Compilador}: GCC 16.1.1 20260430
    \item \textbf{Soporte OpenMP}: OpenMP 5.1 (GCC libgomp)
    \item \textbf{Banderas de compilaci\'on}: \texttt{g++ -O3 -fopenmp -march=native}
    \item \textbf{Tama\~no de datos}: $50\,000\,000 \times 16 = 6.4$ GB
    \item \textbf{Repeticiones}: Una corrida por configuraci\'on
    \item \textbf{M\'etrica principal}: Tiempo total = Fase 1 + Fase 2 + Fase 3
\end{itemize}

\subsubsection{Tabla de Tiempos}

\begin{table}[H]
\centering
\begin{tabular}{crrrrr}
\toprule
\textbf{Hilos} & \textbf{Fase 1 (s)} & \textbf{Fase 2 (s)} & \textbf{Fase 3 (s)} & \textbf{Total (s)} & \textbf{Speedup $S_p$} \\
\midrule
1 & 0.8340 & 0.7800 & 1.0272 & 2.6412 & 1.000 \\
2 & 0.3126 & 0.4102 & 0.7544 & 1.4773 & 1.788 \\
4 & 0.2059 & 0.2282 & 0.8009 & 1.2350 & 2.139 \\
8 & 0.2469 & 0.2350 & 0.8331 & 1.3150 & 2.008 \\
\bottomrule
\end{tabular}
\caption{Tiempos de ejecuci\'on medidos y speedup respecto a 1 hilo.}
\end{table}

\subsubsection{C\'alculo de Speedup y Eficiencia}

\[
S_p = \frac{T_1}{T_p} \qquad E_p = \frac{S_p}{p}
\]

\begin{table}[H]
\centering
\begin{tabular}{crrr}
\toprule
$p$ & $T_p$ (s) & $S_p$ & $E_p$ (\%) \\
\midrule
1 & 2.6412 & 1.000 & 100.0 \\
2 & 1.4773 & 1.788 & 89.4 \\
4 & 1.2350 & 2.139 & 53.5 \\
8 & 1.3150 & 2.008 & 25.1 \\
\bottomrule
\end{tabular}
\caption{Speedup y eficiencia para cada configuraci\'on de hilos.}
\end{table}

\subsubsection{An\'alisis de Resultados}

El speedup m\'aximo alcanzado es de $2.14\times$ con 4 hilos. Con 8 hilos el
rendimiento decrece ligeramente ($2.01\times$). Este comportamiento es
consistente con un problema limitado por ancho de banda de memoria:

\begin{itemize}
    \item Las fases 1 y 2 son de solo lectura y escalan mejor: $S_4 \approx
        3.6\times$ y $S_8 \approx 3.3\times$ en la fase 1, pues los 6.4 GB
        se leen desde DRAM sin contenci\'on de escritura.
    \item La fase 3 (normalizaci\'on) lee y escribe 6.4 GB, saturando el bus
        de memoria. Con 1 hilo la velocidad efectiva es
        $6.4 / 1.0272 \approx 6.2$ GB/s; con 2 hilos sube a $\sim$8.5 GB/s,
        y con 4 u 8 hilos se estabiliza en $\sim$8 GB/s, indicando saturaci\'on
        del ancho de banda DDR4.
    \item El speedup total es menor que el del min-max original debido a las
        tres pasadas completas sobre los datos.
    \item Se detectaron 4\,734\,853 valores at\'ipicos ($|Z| > 3$), que
        representan el 0.62\% de los datos v\'alidos.
\end{itemize}

\subsubsection{Por qu\'e \texttt{critical} dentro del bucle principal ser\'ia una mala decisi\'on}

Si se usara \texttt{\#pragma omp critical} dentro del bucle de normalizaci\'on
para proteger la escritura de cada celda, se introducir\'ia una serializaci\'on
completa: todos los hilos competir\'ian por la misma secci\'on cr\'itica,
anulando cualquier beneficio del paralelismo. El overhead de entrar y salir
de la secci\'on cr\'itica $800$ millones de veces har\'ia la ejecuci\'on mucho
m\'as lenta que la versi\'on secuencial.

En nuestra implementaci\'on, la secci\'on \texttt{critical} se usa solo una vez
por hilo por fase para combinar 16 o 32 valores, lo que representa un overhead
despreciable.

\subsubsection{Error Num\'erico y Orden de Reducci\'on en Punto Flotante}

La suma de punto flotante no es asociativa debido a la precisi\'on finita de
IEEE 754. En las fases 1 y 2, los acumuladores privados se suman en orden
indeterminado dentro de la secci\'on \texttt{critical}, lo que puede producir
diferencias del orden de $10^{-6}$ en $\mu_j$ y $\sigma_j$ entre ejecuciones.
Este error es despreciable para el an\'alisis de datos y no afecta la
clasificaci\'on de at\'ipicos.

\subsection{Conclusiones del Ejercicio 1}

\begin{enumerate}
    \item El problema es \emph{memory-bound}. El speedup est\'a limitado por
        el ancho de banda de memoria DDR4, no por la capacidad de c\'omputo.
    \item El Z-Score requiere 3 pasadas completas sobre los datos, incrementando
        el tr\'afico de memoria en $\sim$33\% respecto a min-max.
    \item OpenMP permite paralelizar eficientemente usando arreglos privados
        por hilo y secciones \texttt{critical} para combinar.
    \item La estrategia fallar\'ia en sistemas con menos RAM que el tama\~no
        del dataset, requiriendo procesamiento por bloques con lectura a disco.
\end{enumerate}

% ============================================================
\section{Ejercicio 2: Pipeline Concurrente de Procesamiento de Logs en Python}
% ============================================================

\subsection{Diagn\'ostico y Selecci\'on de Estrategia}

Durante la evaluaci\'on individual (Parte I), cada integrante propuso una
estrategia distinta para este problema. Al comparar en grupo, las diferencias
m\'as relevantes giraron en torno a si usar \texttt{threading},
\texttt{multiprocessing} o un enfoque h\'ibrido. La discusi\'on nos llev\'o a
descartar \texttt{threading} como soluci\'on principal y a elegir
\textbf{\texttt{ProcessPoolExecutor} con particionamiento por chunks} por las
siguientes razones:

\begin{itemize}
    \item \textbf{El GIL bloquea las etapas m\'as costosas}: las etapas de
        parseo, limpieza y agregaci\'on son CPU-bound y ejecutan Python puro.
        El GIL impide que dos hilos ejecuten bytecode Python simult\'aneamente,
        por lo que \texttt{threading} no introduce paralelismo real en estas
        fases.
    \item \textbf{La lectura puede ser secuencial}: la descompresi\'on con
        \texttt{gzip} libera el GIL parcialmente (la rutina de inflado est\'a
        en C), por lo que \texttt{threading} ayudar\'ia en la lectura, pero
        ese beneficio es menor comparado con el costo del parseo.
    \item \textbf{El merge es seguro sin locks}: al usar el patr\'on
        \emph{map-reduce} (cada worker retorna un diccionario parcial y el
        proceso principal los combina), se evitan condiciones de carrera sin
        ning\'un mecanismo de sincronizaci\'on expl\'icito. Esto es equivalente
        a \texttt{reduction} en OpenMP.
\end{itemize}

Descartamos \texttt{asyncio} porque los logs est\'an en disco local, no en red.

\subsection{Diagn\'ostico del Problema}

Se requiere procesar 2 GB de logs comprimidos en formato JSONL con los campos:
\texttt{user\_id}, \texttt{timestamp}, \texttt{endpoint}, \texttt{latency\_ms}
y \texttt{status\_code}. El objetivo es agregar m\'etricas por ventanas
temporales de 10 minutos.

\begin{table}[H]
\centering
\begin{tabular}{lll}
\toprule
\textbf{Etapa} & \textbf{Tipo de carga} & \textbf{Raz\'on} \\
\midrule
Lectura y descompresi\'on & I/O-bound & Acceso a disco; \texttt{gzip} usa C interno \\
Parseo JSON              & CPU-bound & \texttt{json.loads()} por l\'inea, Python puro \\
Limpieza y validaci\'on  & CPU-bound & Verificaciones de campos, conversiones de tipo \\
Agregaci\'on por ventanas & CPU-bound & Agrupaciones, sumas, conteos en diccionarios \\
Escritura de resultados  & I/O-bound & Salida peque\~na, no es cuello de botella \\
\bottomrule
\end{tabular}
\caption{Tipo de carga por etapa del pipeline.}
\end{table}

\subsection{Dataset Sint\'etico}

Se gener\'o un dataset sint\'etico de aproximadamente 2 GB comprimidos
($\sim$50 millones de l\'ineas) usando NumPy:

\begin{lstlisting}[style=python, caption={Generaci\'on de dataset sint\'etico.}]
import gzip, json, numpy as np
from datetime import datetime, timedelta

def generate_logs(n=50_000_000, filename="logs.jsonl.gz"):
    users     = [f"user_{i}" for i in range(1000)]
    endpoints = ["/api/search", "/api/checkout", "/api/login", "/api/data"]
    codes     = [200, 200, 200, 404, 500]
    base      = datetime(2024, 1, 1)
    batch     = 100_000
    with gzip.open(filename, "wt") as f:
        for start in range(0, n, batch):
            size = min(batch, n - start)
            u   = np.random.randint(0, 1000, size)
            e   = np.random.randint(0, 4, size)
            c   = np.random.choice(codes, size)
            lat = np.round(np.random.normal(120, 40, size), 2)
            ts_offsets = np.arange(start, start + size) * 0.1
            for i in range(size):
                ts = (base + timedelta(
                    seconds=float(ts_offsets[i]))).isoformat()
                f.write(json.dumps({
                    "user_id":     users[u[i]],
                    "timestamp":   ts,
                    "endpoint":    endpoints[e[i]],
                    "latency_ms":  float(lat[i]),
                    "status_code": int(c[i])
                }) + "\n")
\end{lstlisting}

\textbf{Nota sobre escala}: el benchmark se ejecut\'o sobre 5 millones de
l\'ineas ($\sim$200 MB comprimidos) y los resultados se extrapolan linealmente
a 50 millones, lo cual es v\'alido porque el trabajo por l\'inea es homog\'eneo.

\subsection{Tabla Comparativa de Estrategias}

\begin{table}[H]
\centering
\small
\begin{tabular}{p{2.8cm}p{2.2cm}p{2.2cm}p{2.2cm}p{2.2cm}}
\toprule
\textbf{Criterio} & \textbf{Secuencial} & \textbf{threading /} \newline \textbf{ThreadPool} & \textbf{multiprocessing /} \newline \textbf{ProcessPool} & \textbf{H\'ibrido} \newline \textbf{(asyncio+MP)} \\
\midrule
Paralelismo CPU real        & No & No (GIL) & S\'i & S\'i \\
Paralelismo I/O             & No & S\'i      & S\'i & S\'i \\
Overhead                    & M\'inimo & Bajo & Alto (fork+pickle) & Medio \\
Complejidad c\'odigo        & Baja & Baja & Media & Alta \\
\'Util parseo CPU-bound     & No & No & S\'i & S\'i \\
\'Util lectura I/O-bound    & No & S\'i & S\'i & S\'i \\
\textbf{Recomendado}        & Baseline & No & \textbf{S\'i} & Solo logs en red \\
\bottomrule
\end{tabular}
\caption{Comparaci\'on de estrategias de concurrencia.}
\end{table}

\subsection{Impacto del GIL}

\begin{itemize}
    \item \textbf{Lectura y descompresi\'on}: \texttt{zlib} libera el GIL
        durante el inflado. \texttt{threading} s\'i paralela esta etapa
        parcialmente.
    \item \textbf{Parseo JSON}: \texttt{json.loads()} est\'a en C pero crear
        objetos Python requiere el GIL. Beneficio marginal con
        \texttt{threading}.
    \item \textbf{Limpieza y agregaci\'on}: Python puro. El GIL bloquea
        completamente. Solo \texttt{multiprocessing} es efectivo.
\end{itemize}

\subsection{Estrategia Seleccionada}

\subsubsection{Dise\~no de Particionamiento}

Se particiona el archivo en chunks de 50.000 l\'ineas. Cada chunk es una lista
de strings enviada a un worker. El worker retorna un diccionario parcial con
clave $(\texttt{endpoint},\; \texttt{ventana\_10min})$ y valor
$\{\texttt{count},\; \texttt{total\_latency},\; \texttt{errors}\}$.
El proceso principal combina los diccionarios con una reducci\'on secuencial
sin locks.

\subsubsection{Pseudoc\'odigo}

\begin{lstlisting}[style=python, caption={Pipeline con ProcessPoolExecutor.}]
from concurrent.futures import ProcessPoolExecutor
from collections import defaultdict
import gzip, json

def read_chunks(filepath, chunk_size=50_000):
    with gzip.open(filepath, "rt") as f:
        chunk = []
        for line in f:
            chunk.append(line)
            if len(chunk) >= chunk_size:
                yield chunk
                chunk = []
        if chunk:
            yield chunk

def process_chunk(lines):
    partial = defaultdict(lambda: {"count": 0,
                                   "total_latency": 0.0,
                                   "errors": 0})
    for line in lines:
        try:
            r = json.loads(line)
            if r["latency_ms"] < 0:
                continue
            key = (r["endpoint"], r["timestamp"][:15] + "0:00")
            partial[key]["count"]         += 1
            partial[key]["total_latency"] += r["latency_ms"]
            if r["status_code"] >= 400:
                partial[key]["errors"]    += 1
        except (json.JSONDecodeError, KeyError):
            continue
    return dict(partial)

def merge_results(partials):
    final = defaultdict(lambda: {"count": 0,
                                 "total_latency": 0.0,
                                 "errors": 0})
    for partial in partials:
        for key, vals in partial.items():
            final[key]["count"]         += vals["count"]
            final[key]["total_latency"] += vals["total_latency"]
            final[key]["errors"]        += vals["errors"]
    return final

def run_pipeline(filepath, n_workers=4):
    import time
    t0       = time.perf_counter()
    chunks   = list(read_chunks(filepath))
    with ProcessPoolExecutor(max_workers=n_workers) as pool:
        partials = list(pool.map(process_chunk, chunks))
    final    = merge_results(partials)
    return final, time.perf_counter() - t0
\end{lstlisting}

El merge no requiere locks porque todos los workers terminan antes de que
\texttt{merge\_results} comience (barrera impl\'icita al salir del bloque
\texttt{with}). Es el equivalente directo de \texttt{reduction(+:...)} en
OpenMP.

\subsection{Identificaci\'on de Overhead}

\begin{table}[H]
\centering
\begin{tabular}{p{3cm}p{4cm}p{3.5cm}p{3cm}}
\toprule
\textbf{Fuente} & \textbf{D\'onde ocurre} & \textbf{Magnitud} & \textbf{Mitigaci\'on} \\
\midrule
Serializaci\'on pickle & Env\'io de chunks y recepci\'on de dicts & Alta si chunks $<$ 10K l\'ineas & Chunks de 50K ($\approx$5 MB) \\
Fork de procesos & Creaci\'on del pool & Fijo $\sim$100--300 ms & Reutilizar el pool \\
Copia IPC & Worker $\to$ proceso principal & Proporcional al dict parcial & Devolver solo agregaciones \\
Merge secuencial & Proceso principal & $<$ 1\% del tiempo total & No paralelizar \\
\bottomrule
\end{tabular}
\caption{Fuentes de overhead y mitigaciones.}
\end{table}

\subsection{Benchmark y Resultados}

\subsubsection{Registro del Entorno de Ejecuci\'on}

\begin{itemize}
    \item \textbf{Sistema operativo}: Windows 11
    \item \textbf{CPU}: AMD Ryzen 5 5500, 6 n\'ucleos f\'isicos, 12 hilos l\'ogicos
    \item \textbf{RAM}: 32 GB 
    \item \textbf{Versi\'on Python}: 3.12.3
    \item \textbf{Tama\~no de datos medido}: 5.000.000 l\'ineas ($\approx$200 MB comprimidos)
    \item \textbf{Tama\~no de datos objetivo}: 50.000.000 l\'ineas ($\approx$2 GB comprimidos)
    \item \textbf{Repeticiones}: 5 por configuraci\'on; 2 warm-up descartadas
    \item \textbf{Chunk size}: 50.000 l\'ineas
    \item \textbf{M\'etrica}: tiempo wall-clock total (lectura + procesamiento + merge)
\end{itemize}

\subsubsection{Tabla de Tiempos Medidos (5M l\'ineas)}

\begin{table}[H]
\centering
\begin{tabular}{crrrr}
\toprule
\textbf{Estrategia} & \textbf{Workers} & \textbf{Tiempo (s)} & \textbf{Speedup $S_p$} & \textbf{Eficiencia $E_p$ (\%)} \\
\midrule
Secuencial               & 1 & 15.57 & 1.00 & 100.0 \\
\texttt{threading}       & 4 & 15.18 & 1.03 & 25.7  \\
\texttt{multiprocessing} & 2 & 8.31  & 1.87 & 93.7  \\
\texttt{multiprocessing} & 4 & 4.81  & 3.24 & 80.9  \\
\texttt{multiprocessing} & 8 & 3.44  & 4.53 & 56.7  \\
\bottomrule
\end{tabular}
\caption{Tiempos medidos sobre 5M l\'ineas.}
\end{table}
\[
S_p = \frac{T_1}{T_p} \qquad E_p = \frac{S_p}{p}
\]

Con 4 workers: $S_4 = 15.57/4.81 = 3.24$, $E_4 = 80.9\%$. \quad
Con 8 workers: $S_8 = 15.57/3.44 = 4.53$, $E_8 = 56.7\%$.

\subsubsection{An\'alisis de Resultados}

\begin{itemize}
    \item \textbf{\texttt{threading} no escala}: speedup de $1.03\times$
        confirma que el GIL anula el paralelismo CPU-bound.
    \item \textbf{\texttt{multiprocessing} escala bien}: eficiencia de 93.7\%
        con 2 workers y 80.9\% con 4 workers, mostrando overhead moderado
        de pickle.
    \item \textbf{Saturaci\'on a 8 workers}: con 8 workers se supera el
        n\'umero de n\'ucleos f\'isicos (6), por lo que los workers
        adicionales operan sobre hilos l\'ogicos via hyperthreading,
        explicando la menor eficiencia (56.7\%).
    \item \textbf{Extrapolaci\'on lineal v\'alida}: el n\'umero de grupos de
        agregaci\'on es finito ($\approx$35.000 grupos), no crece
        superlinealmente.
\end{itemize}

\subsection{Riesgos y Limitaciones}

\begin{table}[H]
\centering
\begin{tabular}{p{3cm}p{5cm}p{5cm}}
\toprule
\textbf{Riesgo} & \textbf{Descripci\'on} & \textbf{Condici\'on de fallo} \\
\midrule
Cuello de botella en lectura & Lectura secuencial limita el throughput & Disco lento (HDD) o archivo en red \\
OOM en merge & Dict final crece si hay muchas claves \'unicas & Ventanas $<$ 1 min o millones de endpoints \\
Overhead domina & Pickle supera el beneficio con chunks peque\~nos & chunk\_size $<$ 10K con muchos workers \\
\texttt{asyncio} ineficaz & No ayuda con I/O de disco local & Solo si logs vienen de S3 o HTTP \\
\bottomrule
\end{tabular}
\caption{Riesgos identificados y condiciones de fallo.}
\end{table}

\subsection{Conclusiones del Ejercicio 2}

\begin{enumerate}
    \item No existe una herramienta universal: \texttt{multiprocessing} es la
        \'unica opci\'on efectiva en Python para pipelines CPU-bound con logs
        en disco local.
    \item El GIL hace que \texttt{threading} sea ineficaz en las etapas
        dominantes; el speedup real es $\approx 1.0$.
    \item El patr\'on map-reduce elimina la necesidad de locks por dise\~no,
        equivalente a \texttt{reduction} en OpenMP.
    \item El speedup m\'aximo real es $\approx 4\times$ antes de saturar la
        lectura. Dividir el dataset en m\'ultiples archivos permitir\'ia
        paralelizar tambi\'en la lectura.
\end{enumerate}

% ============================================================
\section{Ejercicio 3: C\'alculo de Similitud por Bloques para Embeddings}
% ============================================================

\subsection{Diagn\'ostico y Selecci\'on de Estrategia}

El problema consiste en calcular similitudes entre $20\,000$ embeddings de
dimensi\'on $128$, normalizados a norma unitaria, y obtener para cada vector
sus $10$ vecinos m\'as similares. La similitud utilizada corresponde a
similitud coseno. Como todos los vectores est\'an normalizados a norma $1$, la
similitud coseno se reduce al producto punto:

\[
\operatorname{sim}(x_i,x_j)=
\frac{x_i \cdot x_j}{\|x_i\|_2\|x_j\|_2}
= x_i \cdot x_j
\]

La estrategia ingenua consiste en construir la matriz completa de similitud
$S = XX^T$, donde $X \in \mathbb{R}^{20000 \times 128}$ y
$S \in \mathbb{R}^{20000 \times 20000}$. Sin embargo, esta estrategia almacena
$400\,000\,000$ valores aunque el objetivo final solo requiere $10$ valores por
fila. Por lo tanto, se descarta la materializaci\'on completa de la matriz y se
selecciona una estrategia exacta basada en procesamiento por bloques y
actualizaci\'on incremental del top-$k$.

La soluci\'on implementada utiliza Python con \texttt{NumPy} y
\texttt{multiprocessing}. No se usa \texttt{g++} ni OpenMP en esta parte, porque
el ejercicio exige una estrategia de bloques, memoria, top-$k$ y m\'etricas de
rendimiento, pero no obliga a implementar OpenMP espec\'ificamente para este
tercer ejercicio. Para evitar la limitaci\'on del GIL de Python, se utilizan
procesos mediante \texttt{ProcessPoolExecutor}. La matriz de embeddings se
comparte entre procesos con \texttt{shared\_memory}, evitando copiar el dataset
completo a cada worker.

\subsection{Diagn\'ostico del Problema}

Sea $N=20\,000$ el n\'umero de vectores y $D=128$ la dimensi\'on de cada
embedding. Si se comparan todos los vectores contra todos los vectores, el
n\'umero de comparaciones es:

\[
N(N-1) = 20\,000 \times 19\,999 = 399\,980\,000
\]

Si se considera tambi\'en la diagonal, la matriz completa tiene:

\[
N^2 = 20\,000^2 = 400\,000\,000
\]

entradas. Cada comparaci\'on requiere aproximadamente $D$ multiplicaciones y
sumas, por lo que el costo computacional es:

\[
O(N^2D)
\]

Para este caso:

\[
400\,000\,000 \times 128 = 51\,200\,000\,000
\]

operaciones elementales de multiplicaci\'on-acumulaci\'on aproximadas. Por lo
tanto, el problema no se resuelve simplemente paralelizando un doble bucle. La
dificultad real est\'a en controlar memoria, localidad de acceso, top-$k$
parcial, reducci\'on de candidatos y costo cuadr\'atico.

\subsection{Dataset Utilizado}

Se gener\'o un dataset sint\'etico compuesto por $20\,000$ vectores de
dimensi\'on $128$. Los valores se generaron con distribuci\'on normal est\'andar
y posteriormente cada fila fue normalizada por su norma euclidiana:

\[
x_i \leftarrow \frac{x_i}{\|x_i\|_2}
\]

Esto garantiza que:

\[
\|x_i\|_2 = 1
\]

La generaci\'on sint\'etica es adecuada para este ejercicio porque permite
controlar exactamente el tama\~no del problema, cumplir con el formato exigido
y reproducir los resultados usando una semilla fija.

\begin{lstlisting}[caption={Generaci\'on de embeddings sint\'eticos normalizados.}]
def generar_embeddings(N, D, seed=42):
    rng = np.random.default_rng(seed)
    X = rng.normal(0, 1, size=(N, D)).astype(np.float32)

    normas = np.linalg.norm(X, axis=1, keepdims=True)
    normas[normas == 0] = 1.0

    X = X / normas
    return X.astype(np.float32)
\end{lstlisting}

\subsection{Estimaci\'on de Memoria}

\subsubsection{Matriz completa de similitud}

La matriz completa de similitud tendr\'ia $400\,000\,000$ elementos. En
precisi\'on \texttt{float32}, cada elemento ocupa $4$ bytes:

\[
400\,000\,000 \times 4 = 1\,600\,000\,000 \text{ bytes}
\]

Esto equivale aproximadamente a:

\[
1.49 \text{ GiB} \approx 1.6 \text{ GB}
\]

En \texttt{float64}, cada elemento ocupa $8$ bytes:

\[
400\,000\,000 \times 8 = 3\,200\,000\,000 \text{ bytes}
\]

Esto equivale aproximadamente a:

\[
2.98 \text{ GiB} \approx 3.2 \text{ GB}
\]

Estas cifras solo consideran la matriz de similitud. En una ejecuci\'on real
tambi\'en se requieren memoria para el dataset, \'indices, estructuras de
ordenamiento y temporales internos. Adem\'as, almacenar toda la matriz es
ineficiente porque solo se conservan $10$ vecinos por cada vector.

\subsubsection{Memoria de la matriz de embeddings}

La matriz original $X$ en \texttt{float32} ocupa:

\[
20\,000 \times 128 \times 4 = 10\,240\,000 \text{ bytes}
\]

\[
\approx 9.77 \text{ MiB}
\]

\subsubsection{Memoria de top-$k$}

Para guardar los $10$ mejores valores y sus \'indices por cada vector:

\[
20\,000 \times 10 \times 4 = 800\,000 \text{ bytes}
\]

para valores, y la misma cantidad para \'indices enteros de 32 bits. En total:

\[
1\,600\,000 \text{ bytes} \approx 1.53 \text{ MiB}
\]

\subsubsection{Memoria temporal por bloques}

En la implementaci\'on usada, las filas se particionan entre workers y las
columnas se procesan en bloques de tama\~no $B=1024$. Cada worker calcula una
submatriz temporal:

\[
S_{\text{temp}} = X_{\text{filas worker}} X_{\text{bloque columnas}}^T
\]

Con $8$ workers, cada worker procesa aproximadamente:

\[
\frac{20\,000}{8}=2\,500
\]

filas. La matriz temporal por worker tiene entonces:

\[
2\,500 \times 1024 = 2\,560\,000
\]

valores \texttt{float32}, lo que equivale aproximadamente a $9.77$ MiB por
worker. La memoria temporal agregada para los $8$ workers se mantiene cercana a
$78$ MiB, sin considerar temporales internos de \texttt{NumPy}. Aun as\'i, este
consumo es muy inferior a mantener la matriz completa de similitud.

\subsection{Comparaci\'on de Estrategias}

\subsubsection{Estrategia A: matriz completa m\'as ordenamiento}

Esta alternativa calcula directamente:

\[
S = XX^T
\]

Luego ordena cada fila de $S$ para obtener los $10$ mayores valores. Su ventaja
principal es la simplicidad conceptual y la posibilidad de usar multiplicaci\'on
matricial altamente optimizada. Sin embargo, presenta tres problemas:

\begin{itemize}
    \item Materializa $400\,000\,000$ similitudes aunque solo se requieren
        $20\,000 \times 10$ resultados finales.
    \item Requiere al menos $1.49$ GiB en \texttt{float32} solo para la matriz
        de similitud.
    \item Ordenar filas completas de $20\,000$ elementos es innecesario cuando
        solo se busca top-$10$.
\end{itemize}

Por estas razones, se descarta como estrategia final.

\subsubsection{Estrategia B: bloques exactos con top-$k$ incremental}

Esta alternativa divide el c\'alculo en bloques. Cada bloque de columnas se
compara contra un subconjunto de filas y se actualiza el top-$10$ de cada vector
sin almacenar todas las similitudes. El bloque temporal se descarta luego de
actualizar los candidatos.

La estrategia mantiene exactitud porque todos los pares son comparados, pero no
materializa la matriz global. Su costo computacional sigue siendo $O(N^2D)$,
pero el consumo de memoria se reduce significativamente. Esta es la estrategia
seleccionada.

\subsubsection{Estrategia C: b\'usqueda aproximada de vecinos}

Una alternativa para escenarios de mayor escala es usar t\'ecnicas aproximadas
como ANN, HNSW o FAISS. Estas reducen el n\'umero de comparaciones y escalan
mejor a millones de vectores. Sin embargo, no garantizan recuperar exactamente
el top-$10$ real, por lo que requieren una m\'etrica de calidad como:

\[
Recall@10 =
\frac{|Top10_{\text{exacto}} \cap Top10_{\text{aproximado}}|}{10}
\]

No se selecciona esta alternativa porque el tama\~no exigido en el enunciado
puede resolverse con una estrategia exacta por bloques.

\subsection{Estrategia Seleccionada}

La estrategia final consiste en:

\begin{enumerate}
    \item Generar o cargar una matriz $X$ de $20\,000 \times 128$ en
        \texttt{float32}.
    \item Normalizar cada vector a norma $1$.
    \item Compartir $X$ entre procesos usando \texttt{shared\_memory}.
    \item Particionar las filas entre workers.
    \item Recorrer las columnas en bloques de tama\~no $B=1024$.
    \item Calcular similitudes parciales con multiplicaci\'on matricial:
        \[
        S_{\text{bloque}} = X_{\text{filas}}X_{\text{columnas}}^T
        \]
    \item Eliminar la diagonal asignando $-\infty$ a la comparaci\'on
        $x_i$ contra s\'i mismo.
    \item Actualizar top-$10$ mediante \texttt{argpartition}, sin ordenar todos
        los candidatos.
    \item Retornar los top-$10$ parciales de cada bloque de filas y construir
        el resultado global.
\end{enumerate}

\subsubsection{Pseudoc\'odigo}

\begin{lstlisting}[caption={Pseudoc\'odigo de la estrategia por bloques.}]
Entrada: X normalizada, N, D, K=10, B=1024, workers

Compartir X en memoria compartida

Dividir filas de X entre workers

Para cada worker en paralelo:
    recibir rango de filas [inicio, fin)
    inicializar top_values con -infinito
    inicializar top_indices con -1

    Para col_start desde 0 hasta N con paso B:
        col_end = min(col_start + B, N)
        Xj = X[col_start:col_end]

        sims = X[inicio:fin] @ Xj.T

        Si una fila se compara consigo misma:
            sims[fila_local, columna_local] = -infinito

        candidatos = concatenar(top_actual, sims)
        indices = concatenar(indices_actuales, indices_del_bloque)

        conservar solo los K mayores con argpartition

    ordenar top-K final de mayor a menor
    retornar top_values y top_indices

Unir resultados parciales de todos los workers
\end{lstlisting}

\subsubsection{Fragmento de C\'odigo Relevante}

\begin{lstlisting}[caption={Worker que procesa un rango de filas y actualiza top-$k$.}]
def worker_topk(args):
    row_start, row_end, N, D, B, K, shm_name = args

    shm = shared_memory.SharedMemory(name=shm_name)
    X = np.ndarray((N, D), dtype=np.float32, buffer=shm.buf)

    row_count = row_end - row_start
    Xi = X[row_start:row_end]

    top_values = np.full((row_count, K), -np.inf, dtype=np.float32)
    top_indices = np.full((row_count, K), -1, dtype=np.int32)

    row_ids = np.arange(row_start, row_end, dtype=np.int32)

    for col_start in range(0, N, B):
        col_end = min(col_start + B, N)
        Xj = X[col_start:col_end]

        sims = Xi @ Xj.T

        mask = (row_ids >= col_start) & (row_ids < col_end)
        if np.any(mask):
            local_rows = np.where(mask)[0]
            local_cols = row_ids[mask] - col_start
            sims[local_rows, local_cols] = -np.inf

        col_indices = np.arange(col_start, col_end, dtype=np.int32)
        col_indices_matrix = np.broadcast_to(col_indices, sims.shape)

        candidatos_values = np.concatenate((top_values, sims), axis=1)
        candidatos_indices = np.concatenate((top_indices, col_indices_matrix), axis=1)

        pos = np.argpartition(candidatos_values, -K, axis=1)[:, -K:]

        top_values = np.take_along_axis(candidatos_values, pos, axis=1)
        top_indices = np.take_along_axis(candidatos_indices, pos, axis=1)

    orden = np.argsort(-top_values, axis=1)
    top_values = np.take_along_axis(top_values, orden, axis=1)
    top_indices = np.take_along_axis(top_indices, orden, axis=1)

    shm.close()
    return row_start, top_values, top_indices
\end{lstlisting}

\subsection{Paralelizaci\'on, Balance de Carga y Localidad de Memoria}

\subsubsection{Paralelizaci\'on}

La paralelizaci\'on se realiza por partici\'on de filas. Cada proceso recibe un
rango de filas distinto y calcula el top-$10$ de esos vectores compar\'andolos
contra todos los bloques de columnas. La matriz $X$ es compartida en modo solo
lectura, por lo que no hay condiciones de carrera sobre los embeddings.

Cada worker mantiene sus propios arreglos locales \texttt{top\_values} y
\texttt{top\_indices}. Como ning\'un worker escribe el top-$k$ de filas ajenas,
no se requiere bloqueo, \texttt{critical}, \texttt{lock} ni sincronizaci\'on fina
durante el c\'alculo principal.

\subsubsection{Balance de carga}

El balance de carga es adecuado porque todos los workers reciben una cantidad
similar de filas. Cada fila debe compararse contra los $N$ vectores, por lo que
el trabajo asignado a cada worker es aproximadamente:

\[
\frac{N}{p} \times N \times D
\]

donde $p$ es el n\'umero de workers. Para $N=20\,000$ y $p=8$, cada worker
procesa aproximadamente $2\,500$ filas contra los $20\,000$ vectores.

\subsubsection{Localidad de memoria}

La localidad de memoria mejora porque las columnas se recorren en bloques de
$1024$ vectores. Cada bloque $X_j$ se reutiliza para calcular similitudes contra
todas las filas asignadas al worker. Adem\'as, al mantener top-$k$ local se
evita escribir una matriz global de $N \times N$ similitudes.

Sin embargo, el algoritmo sigue generando matrices temporales para cada bloque,
por lo que el rendimiento puede quedar limitado por ancho de banda de memoria,
creaci\'on de temporales y operaciones de selecci\'on top-$k$.

\subsection{Benchmark y Resultados}

\subsubsection{Registro del entorno de ejecuci\'on}

\begin{itemize}
    \item \textbf{Sistema operativo}: Windows, ejecuci\'on desde PowerShell.
    \item \textbf{Lenguaje}: Python.
    \item \textbf{Librer\'ias principales}: \texttt{NumPy},
        \texttt{multiprocessing}, \texttt{concurrent.futures}.
    \item \textbf{Archivo ejecutado}: \texttt{Ejercicio\_3.py}.
    \item \textbf{Tama\~no de datos}: $20\,000 \times 128$ embeddings.
    \item \textbf{Tipo de dato}: \texttt{float32}.
    \item \textbf{Top-k}: $K=10$.
    \item \textbf{Tama\~no de bloque}: $B=1024$.
    \item \textbf{Workers evaluados}: $1$, $2$, $4$ y $8$.
    \item \textbf{Repeticiones}: una ejecuci\'on por configuraci\'on.
    \item \textbf{CPU y RAM}: no registradas expl\'icitamente en la captura;
        se recomienda agregarlas desde el equipo usado antes de la entrega
        final.
\end{itemize}

\subsubsection{M\'etricas utilizadas}

Las m\'etricas evaluadas fueron:

\[
S_p = \frac{T_1}{T_p}
\]

\[
E_p = \frac{S_p}{p}
\]

\[
Throughput = \frac{N(N-1)}{T_p}
\]

donde $T_1$ es el tiempo con un worker, $T_p$ es el tiempo con $p$ workers y
$N(N-1)$ corresponde al n\'umero de comparaciones excluyendo la diagonal.

\subsubsection{Tabla de benchmark}

\begin{table}[H]
\centering
\begin{tabular}{crrrr}
\toprule
\textbf{Workers} & \textbf{Tiempo (s)} & \textbf{Speedup} & \textbf{Eficiencia} & \textbf{Throughput comp/s} \\
\midrule
1 & 4.1124 & 1.0000 & 1.0000 & 97\,263\,029.82 \\
2 & 2.3791 & 1.7286 & 0.8643 & 168\,124\,738.32 \\
4 & 1.8001 & 2.2845 & 0.5711 & 222\,200\,717.06 \\
8 & 1.6417 & 2.5049 & 0.3131 & 243\,636\,060.18 \\
\bottomrule
\end{tabular}
\caption{Resultados experimentales del c\'alculo de similitud por bloques.}
\end{table}

\subsubsection{An\'alisis de resultados}

Los resultados muestran una mejora efectiva al aumentar el n\'umero de workers.
Con $1$ worker el tiempo total fue de $4.1124$ segundos, mientras que con $8$
workers se redujo a $1.6417$ segundos. Esto corresponde a un speedup m\'aximo
de:

\[
S_8 = \frac{4.1124}{1.6417} = 2.5049
\]

El throughput tambi\'en mejora desde $97.26$ millones de comparaciones por
segundo con $1$ worker hasta $243.64$ millones de comparaciones por segundo con
$8$ workers.

Sin embargo, la escalabilidad no es lineal. La eficiencia cae desde $0.8643$ con
$2$ workers hasta $0.3131$ con $8$ workers. Esto indica que el algoritmo escala,
pero comienza a saturarse por overhead de paralelizaci\'on y presi\'on sobre la
memoria. Entre las causas probables est\'an:

\begin{itemize}
    \item creaci\'on de procesos;
    \item acceso concurrente a memoria compartida;
    \item matrices temporales \texttt{sims};
    \item operaciones \texttt{concatenate};
    \item selecci\'on parcial con \texttt{argpartition};
    \item copias de resultados parciales desde los workers al proceso principal;
    \item limitaciones del ancho de banda de memoria.
\end{itemize}

La mejor configuraci\'on en tiempo total fue $8$ workers. No obstante, desde el
punto de vista de eficiencia paralela, $2$ workers aprovechan mejor los recursos,
ya que entregan $1.7286\times$ de speedup con una eficiencia de $86.43\%$.

\subsection{M\'etrica de Calidad}

La estrategia implementada es exacta, no aproximada. Todos los vectores son
comparados contra todos los dem\'as vectores, y el top-$10$ se obtiene mediante
selecci\'on parcial. Por lo tanto, no se requiere una m\'etrica de calidad como
\texttt{Recall@10}, que ser\'ia necesaria si se usara una t\'ecnica aproximada.

Aun as\'i, se propone validar la correcci\'on comparando el resultado por bloques
contra una implementaci\'on de referencia con matriz completa para un subconjunto
peque\~no, por ejemplo $N=1000$. Si los \'indices top-$10$ coinciden para cada
fila, se valida que la estrategia por bloques conserva el resultado exacto sin
almacenar toda la matriz.

\subsection{Riesgos, Limitaciones y Condiciones de Falla}

\begin{enumerate}
    \item \textbf{Costo cuadr\'atico}: aunque se reduce el consumo de memoria,
        el algoritmo sigue realizando $O(N^2D)$ operaciones. Para datasets de
        cientos de miles o millones de embeddings, el enfoque exacto puede dejar
        de ser viable.
    \item \textbf{Escalabilidad sublineal}: los resultados muestran que aumentar
        workers no genera speedup proporcional. A partir de cierto punto, el
        overhead y el ancho de banda de memoria dominan.
    \item \textbf{Memoria temporal}: aunque no se almacena la matriz completa,
        cada worker crea matrices temporales. Con bloques demasiado grandes,
        la memoria puede crecer y afectar el rendimiento.
    \item \textbf{Tama\~no de bloque}: bloques peque\~nos aumentan overhead de
        iteraci\'on; bloques grandes aumentan uso de memoria. El valor $B=1024$
        representa un compromiso pr\'actico.
    \item \textbf{Diagonal}: si no se elimina la comparaci\'on $x_i$ contra
        $x_i$, cada vector aparecer\'ia como su propio vecino m\'as similar.
    \item \textbf{Dependencia del hardware}: el rendimiento depende del n\'umero
        real de n\'ucleos, ancho de banda de memoria, implementaci\'on BLAS de
        \texttt{NumPy} y carga del sistema.
    \item \textbf{Aproximaciones futuras}: para vol\'umenes mayores se deber\'ia
        considerar FAISS, HNSW o ANN, midiendo calidad mediante \texttt{Recall@10}
        contra una referencia exacta.
\end{enumerate}

\subsection{Conclusiones}

\begin{enumerate}
    \item La matriz completa de similitud requiere aproximadamente $1.49$ GiB
        en \texttt{float32} y $2.98$ GiB en \texttt{float64}, sin contar
        temporales ni estructuras adicionales.
    \item La estrategia por bloques evita materializar la matriz completa y
        mantiene solo el top-$10$ por vector.
    \item La partici\'on por filas evita condiciones de carrera porque cada
        worker actualiza \'unicamente sus propios resultados.
    \item El mejor tiempo observado fue $1.6417$ segundos con $8$ workers,
        obteniendo un speedup de $2.5049\times$.
    \item La eficiencia cae al aumentar workers, lo que evidencia saturaci\'on
        por overhead y memoria. Esto es esperable en un problema con alto volumen
        de datos y selecci\'on top-$k$ incremental.
    \item La estrategia seleccionada es exacta; no requiere una m\'etrica de
        calidad aproximada, aunque se recomienda validar contra una matriz
        completa en un subconjunto peque\~no.
\end{enumerate}

\subsection{Declaraci\'on de Uso de Herramientas}

Se utiliz\'o apoyo de IA generativa para estructurar el informe, revisar la
justificaci\'on t\'ecnica y adaptar la explicaci\'on a los criterios de la
evaluaci\'on. El c\'odigo fue ejecutado y validado experimentalmente por el
grupo, y los resultados incluidos corresponden a la corrida mostrada para el
archivo \texttt{Ejercicio\_3.py}.

\subsection{Referencias}

\begin{enumerate}
    \item Grama, A., Gupta, A., Karypis, G., \& Kumar, V. (2003).
    \textit{Introduction to Parallel Computing}. Addison-Wesley.
    \item Pacheco, P. (2011). \textit{An Introduction to Parallel Programming}.
    Morgan Kaufmann.
    \item NumPy Developers. \textit{NumPy Documentation}. Disponible en:
    \url{https://numpy.org/doc/}
    \item Python Software Foundation. \textit{multiprocessing documentation}.
    Disponible en: \url{https://docs.python.org/3/library/multiprocessing.html}
    \item INF8090 -- Computaci\'on Paralela y Distribuida. Enunciado de
    Evaluaci\'on 1, Segunda Parte, Ejercicio 3.
\end{enumerate}

\end{document}
